% Generated by simulations/python/experiments.py -- do not edit by hand
\begin{table}[htbp]\centering\footnotesize
\setlength{\tabcolsep}{4pt}
\caption{Share of elections a coalition can manipulate}\label{tab:manipulability}
\sbox{\tabbox}{\begin{tabular}{l>{\raggedleft\arraybackslash}p{0.089\linewidth}>{\raggedleft\arraybackslash}p{0.089\linewidth}>{\raggedleft\arraybackslash}p{0.100\linewidth}>{\raggedleft\arraybackslash}p{0.100\linewidth}>{\raggedleft\arraybackslash}p{0.100\linewidth}>{\raggedleft\arraybackslash}p{0.100\linewidth}}
\toprule
Rule & Spatial: all (\%) & Spatial: 10\% (\%) & Polarised: all (\%) & Polarised: 10\% (\%) & Impartial: all (\%) & Impartial: 10\% (\%) \\
\midrule
Plurality & 94.2 & 53.8 & 98.6 & 45.0 & 100.0 & 98.0 \\
IRV & 47.8 & 36.6 & 76.2 & 44.6 & 53.4 & 51.6 \\
Borda & 97.8 & 79.0 & 96.8 & 81.6 & 100.0 & 95.2 \\
Copeland & 96.2 & 19.8 & 96.8 & 31.2 & 100.0 & 88.4 \\
Minimax & 63.2 & 4.8 & 53.8 & 7.0 & 100.0 & 70.8 \\
Schulze & 63.2 & 4.8 & 53.8 & 6.8 & 100.0 & 71.2 \\
Ranked Pairs & 71.6 & 4.4 & 69.8 & 6.2 & 100.0 & 64.8 \\
Split Cycle & 73.0 & 5.0 & 72.2 & 7.2 & 100.0 & 72.2 \\
\midrule
Approval & 99.6 & 84.6 & 100.0 & 68.8 & 100.0 & 90.6 \\
Majority Judgment & 98.8 & 79.0 & 99.4 & 79.8 & 100.0 & 86.6 \\
\midrule
Fallback & 83.6 & 97.4 & 86.2 & 98.8 & 100.0 & 90.6 \\
Fallback (impasse) & 5.8 & 0.0 & 15.8 & 0.0 & 0.0 & 0.0 \\
Majoritarian compromise & 77.2 & 33.0 & 81.2 & 39.4 & 100.0 & 81.6 \\
Sequential veto & 76.8 & 76.8 & 63.2 & 63.2 & 69.2 & 69.2 \\
\midrule
Unanimity & 0.0 & 0.0 & 0.8 & 0.8 & 0.0 & 0.0 \\
PVC + Schulze & 98.4 & 5.8 & 87.6 & 7.0 & 100.0 & 95.4 \\
QVC, any veto (1/n) & 99.6 & 84.6 & 100.0 & 68.8 & 100.0 & 90.6 \\
QVC, 10\% & 99.6 & 84.6 & 100.0 & 68.8 & 100.0 & 90.6 \\
QVC, 25\% & 99.6 & 51.8 & 100.0 & 56.8 & 100.0 & 90.6 \\
QVC, 33\% & 99.4 & 11.0 & 100.0 & 26.6 & 100.0 & 81.8 \\
QVC, 50\% & 72.2 & 4.8 & 83.6 & 7.0 & 100.0 & 71.2 \\
\bottomrule
\end{tabular}}
\usebox{\tabbox}
\tabnotes{500 elections per culture, 101 voters, 8 options. An election counts as manipulable if, for some option $y$, voters who prefer $y$ to the winner can elect $y$ by one of three joint strategies: rank $y$ first and the winner last, vetoing the winner; the same, vetoing every option but $y$; or rank the winner last and veto it. All: every voter who prefers $y$ to the winner joins; 10\%: at most 10 of them. Other strategies may succeed where these fail: the shares are lower bounds. Monte Carlo standard errors at most 2.2 points.}
{Monte Carlo simulation, cultures of \texttt{cultures.py}, rules of \texttt{voting.py}; \texttt{experiments.py}.}
\end{table}
